\section{Définitions}
\label{definitions}

\begin{longtblr}[
 label = none,
 ]{
 colspec = {Q[l,wd=4cm]X[l]},
 width = \linewidth,
 rowsep = 4pt,
 row{1} = {font=\bfseries},
 column{1} = {font=\bfseries},
 column{1} = {m},
 }
 Terme                            & Définition                                                                                                                                                                                                                  \\
 \hline
 FCEG                             & Abréviation de la Fédération canadienne étudiante en génie                                                                                                                                                                  \\
 Statuts constitutifs             & Les statuts constitutifs originaux ou mis à jour ou les statuts de modification, de fusion, de prorogation, de réorganisation, d'arrangement ou de reconstitution de la FCEG.                                               \\
 Résolution                       & Une motion présentée comme proposition et approuvée par un vote majoritaire, ou selon le seuil de vote applicable si spécifié autrement                                                                                     \\
 Résolution spéciale              & Une résolution approuvée par un vote majoritaire des deux tiers (2/3)                                                                                                                                                       \\
 Membre, Association membre       & Une association étudiante de génie reconnue représentant la population étudiante d'un établissement d'enseignement postsecondaire canadien qui a été admise comme membre de la FCEG conformément à la présente constitution \\
 Assemblée générale               & Une réunion des membres                                                                                                                                                                                                     \\
 DGE                              & La direction générale des élections (DGE) est la personne responsable et autorité ultime sur les questions électorales.                                                                                                     \\
 Assemblée de mise en candidature & L'assemblée générale tenue avant l'assemblée générale annuelle aux fins de sélectionner la liste des personnes candidates pour pourvoir les postes vacants à venir.                                                         \\
 Liste de candidatures            & La liste des personnes candidates à l'équipe de direction sélectionnée par les membres lors de l'assemblée de mise en candidature pour présentation lors de l'assemblée générale annuelle.                                  \\
 Exécutif national                & La présidence et les vice-présidences                                                                                                                                                                                       \\
 Équipe de direction              & Le CA et l'exécutif national                                                                                                                                                                                                \\
\end{longtblr}
